\documentclass[11pt,a4paper]{article}
\usepackage[T1]{fontenc}
\usepackage[utf8]{inputenc}
\usepackage{amsmath,amssymb,amsthm}
\usepackage[margin=1in]{geometry}
\usepackage[colorlinks=true,linkcolor=blue,urlcolor=blue]{hyperref}
\theoremstyle{plain}
\newtheorem{theorem}{Theorem}
\newtheorem{lemma}[theorem]{Lemma}
\newtheorem{proposition}[theorem]{Proposition}
\newtheorem{corollary}[theorem]{Corollary}
\theoremstyle{definition}
\newtheorem{remark}[theorem]{Remark}

\title{The empirical recurrence for OEIS A204719, proved by transfer matrix}
\author{Adrian Perez Fontelles\\ \small Independent researcher}
\date{2 September 2026}
\begin{document}
\maketitle

\begin{abstract}
OEIS A204719 counts arrays over $\{0,\dots,1\}$ under a condition on the permanent of every
$2\times2$ subblock, and carries an empirical recurrence of order $59$ contributed by
R.~H.~Hardin. It is true, and it is decidable rather than empirical. The condition is local to
a bounded window of consecutive rows, so the count is a walk count on $S=450$ states and is
$C$-finite; whether the conjectured recurrence annihilates it is then settled by a finite exact
computation, with the Cayley--Hamilton theorem bounding the work.
\end{abstract}

\noindent\small 2020 Mathematics Subject Classification. 05A15, 05B45, 15B36.\normalsize

\section{The conjecture}

OEIS A204719 is ``Number of (n+1)X5 0..1 arrays with the permanents of 2X2 subblocks nondecreasing rightwards and downwards.''. It has offset $1$ and begins
\[
450, 5821, 67518, 727255, 7549845, 75805761, 748847269, 7296210120,\ \dots
\]
The entry states, as a conjecture contributed by its author and never marked settled:
\begin{quote}
Empirical: a(n) = 16*a(n-1) +65*a(n-2) -2014*a(n-3) +243*a(n-4) +116940*a(n-5) -166848*a(n-6) -4177206*a(n-7) +8311505*a(n-8) +103223574*a(n-9) -231273807*a(n-10) -1879943016*a(n-11) +4377822995*a(n-12) +26237776878*a(n-13) -60973013672*a(n-14) -288002292400*a(n-15) +652256570190*a(n-16) +2531735911304*a(n-17) -5507161089400*a(n-18) -18057043358940*a(n-19) +37390631980226*a(n-20) +105483977994876*a(n-21) -206866717988958*a(n-22) -508148160120256*a(n-23) +941620739450405*a(n-24) +2028117686977836*a(n-25) -3550538779616463*a(n-26) -6725933659087278*a(n-27) +11142324614379039*a(n-28) +18557756920384556*a(n-29) -29184370788902344*a(n-30) -42588868781561294*a(n-31) +63873907719689629*a(n-32) +81144901475551854*a(n-33) -116762342636840643*a(n-34) -127899358991420664*a(n-35) +177901630984905193*a(n-36) +165841212216673590*a(n-37) -225045078669344698*a(n-38) -175476958619485124*a(n-39) +234962339938048848*a(n-40) +149787714063731872*a(n-41) -200783588357815520*a(n-42) -101473173031444224*a(n-43) +138829814174390080*a(n-44) +53244993251106048*a(n-45) -76474039873733952*a(n-46) -20808331640441600*a(n-47) +32854281566477504*a(n-48) +5625619348617344*a(n-49) -10684672575669632*a(n-50) -865167444433664*a(n-51) +2517526213069824*a(n-52) +3268384448512*a(n-53) -400848961507328*a(n-54) +27182934855680*a(n-55) +38006887096320*a(n-56) -4794757939200*a(n-57) -1584385228800*a(n-58) +273678336000*a(n-59)
\end{quote}
As of the ``Last modified'' line on the live entry (Oct 03 2025 11:25:57, revision 9) this is still
recorded as empirical, and nothing on the entry records it as proved.

\section{What is being counted}

Write $\sigma$ for the permanent of a $2\times2$ subblock
$\begin{pmatrix}p&q\\r&s\end{pmatrix}$, that is $\sigma=ps+qr$. The subblocks of an
$(n+1)\times5$ array over $\{0,\dots,1\}$ form a grid of $n$ rows and $4$
columns, and $\sigma$ colours that grid.

The condition relates NEIGHBOURING subblocks: the values of two subblocks adjacent in the grid,
horizontally or vertically, must not decrease. Rightwards and downwards are the two directions named.

\section{The count is a walk count}

A subblock's value needs two consecutive array rows, and the vertical
comparison needs the two subblock rows that three consecutive array rows produce. So the state
is the PAIR of consecutive array rows: the horizontal comparisons decide which pairs are
admissible at all, and the vertical ones decide which pair may follow which. A step appends one
array row, turning $(r,s)$ into $(s,t)$ and settling every comparison between the two subblock
rows at once. An $(n+1)$-row array is a walk of $n-1$ steps, so
\[
a(n)\;=\;\iota^{\!\top}M^{\,n-1}\tau,\qquad\iota=\tau=(1,\dots,1)^{\!\top},
\]
on the $S=450$ admissible pairs.

At $n=1$ the array has two rows, so the grid is a single row of $4$
subblocks with no vertical neighbours, so only the horizontal comparisons apply. The
entry's first term is $450$, and the model reproduces it along with every later term; that
is the cheapest check that the grid has been read with the right shape.

Over the alphabet $\{0,\dots,1\}$ the permanent of a $2\times2$ subblock takes values in
$[0,2]$, so the colouring of the grid it induces has at most $3$ values.
That bound is what keeps the construction below finite in the only place where the value of
$\sigma$, rather than a comparison between two of them, has to be carried.
In particular $a$ is $C$-finite. The alignment of the walk index with the entry's own $n$ was
checked against its published terms rather than assumed.

\section{The proof}

\begin{lemma}\label{lem:crit}
Let $q(t)=t^{59}-\sum_i c_it^{59-i}$ be the characteristic polynomial of the
conjectured recurrence. The residual $a(n)-\sum_ic_ia(n-i)$ equals
$\iota^{\!\top}M^{\,j}q(M)\tau$ for the corresponding walk index $j$, so the recurrence
holds from some point on if and only if $\iota^{\!\top}M^{j}q(M)\tau=0$ for all large $j$,
and it is enough to examine $j<S$.
\end{lemma}

\begin{proof}
Factor $M^{j}$ out of $M^{j+59}-\sum_ic_iM^{j+59-i}$. For the bound, $M^{S}$
is an integer combination of $I,M,\dots,M^{S-1}$ by Cayley--Hamilton, so the numbers
$u_j=\iota^{\!\top}M^{j}q(M)\tau$ obey the monic recurrence given by the characteristic
polynomial of $M$; $S$ consecutive zeros force every later one, and the last nonzero $u_j$
pins the threshold exactly.
\end{proof}

\begin{theorem}
\[
a(n)\;=\;16\,a(n-1) + 65\,a(n-2) - 2014\,a(n-3) + 243\,a(n-4) + 116940\,a(n-5) - 166848\,a(n-6) - 4177206\,a(n-7) + 8311505\,a(n-8) + 103223574\,a(n-9) - 231273807\,a(n-10) - 1879943016\,a(n-11) + 4377822995\,a(n-12) + 26237776878\,a(n-13) - 60973013672\,a(n-14) - 288002292400\,a(n-15) + 652256570190\,a(n-16) + 2531735911304\,a(n-17) - 5507161089400\,a(n-18) - 18057043358940\,a(n-19) + 37390631980226\,a(n-20) + 105483977994876\,a(n-21) - 206866717988958\,a(n-22) - 508148160120256\,a(n-23) + 941620739450405\,a(n-24) + 2028117686977836\,a(n-25) - 3550538779616463\,a(n-26) - 6725933659087278\,a(n-27) + 11142324614379039\,a(n-28) + 18557756920384556\,a(n-29) - 29184370788902344\,a(n-30) - 42588868781561294\,a(n-31) + 63873907719689629\,a(n-32) + 81144901475551854\,a(n-33) - 116762342636840643\,a(n-34) - 127899358991420664\,a(n-35) + 177901630984905193\,a(n-36) + 165841212216673590\,a(n-37) - 225045078669344698\,a(n-38) - 175476958619485124\,a(n-39) + 234962339938048848\,a(n-40) + 149787714063731872\,a(n-41) - 200783588357815520\,a(n-42) - 101473173031444224\,a(n-43) + 138829814174390080\,a(n-44) + 53244993251106048\,a(n-45) - 76474039873733952\,a(n-46) - 20808331640441600\,a(n-47) + 32854281566477504\,a(n-48) + 5625619348617344\,a(n-49) - 10684672575669632\,a(n-50) - 865167444433664\,a(n-51) + 2517526213069824\,a(n-52) + 3268384448512\,a(n-53) - 400848961507328\,a(n-54) + 27182934855680\,a(n-55) + 38006887096320\,a(n-56) - 4794757939200\,a(n-57) - 1584385228800\,a(n-58) + 273678336000\,a(n-59)
\]
for every $n>59$.
\end{theorem}

\begin{proof}
The vector $w=q(M)\tau$ was formed by $59$ matrix--vector products in exact integer
arithmetic and $\iota^{\!\top}M^{j}w$ evaluated for $j=0,1,\dots$, again exactly; those
integers vanish from the index corresponding to $n=59+1$ onwards and the run continued
until the working vector was identically zero, which settles every larger $j$ at once.
\end{proof}

\section{Verification}

Three checks, each independent of the algebra above.

First, the model reproduces all $16$ terms the entry publishes, exactly, in integer
arithmetic. This is what ties the model to the entry: the digraph is built by reading the
entry's English, and a misreading gives different counts.

Second, the conjectured recurrence was evaluated directly on those published terms, with no
matrices involved, and holds wherever the proved range and the data overlap.

Third, the annihilation test of Lemma~\ref{lem:crit} was carried out in exact integer
arithmetic on the whole vector rather than on a sampled prefix.

\begin{thebibliography}{9}
\bibitem{oeis} The OEIS Foundation, \emph{The On-Line Encyclopedia of Integer
Sequences}, \url{https://oeis.org}, 2026. Sequence A204719.
\bibitem{stanley} R.~P.~Stanley, \emph{Enumerative Combinatorics}, Volume 1, 2nd ed.,
Cambridge University Press, 2011. (Transfer-matrix method, Section 4.7.)
\bibitem{minc} H.~Minc, \emph{Permanents}, Addison--Wesley, 1978.
\bibitem{horn} R.~A.~Horn and C.~R.~Johnson, \emph{Matrix Analysis}, 2nd ed., Cambridge
University Press, 2013.
\end{thebibliography}

\end{document}
